\documentclass[11pt]{article}
\usepackage[T1]{fontenc}
\usepackage[utf8]{inputenc}
\usepackage{amsmath,amssymb,amsthm,mathtools,xparse,multicol,enumitem,hyperref}
\usepackage[a4paper,margin=25mm]{geometry}
\newtheorem{theorem}{Theorem}[section]
\theoremstyle{definition}
\newtheorem{example}[theorem]{Exercise}
\newtheorem{examples}[theorem]{Exercises}
\newtheorem{problem}[theorem]{Problem}
\theoremstyle{remark}
\newtheorem{remark}[theorem]{Remark}
\newcommand{\pqty}[1]{\left(#1\right)}
\newcommand{\bqty}[1]{\left[#1\right]}
\DeclarePairedDelimiter{\abs}{\lvert}{\rvert}
\newcommand{\flatfrac}[2]{\frac{#1}{#2}}
\newcommand{\set}[2]{\left\{#1\mid #2\right\}}
\newcommand{\ccintv}[2]{[#1,#2]}
\newcommand{\oointv}[2]{(#1,#2)}
\newcommand{\functiondomain}[3]{#1:#2\to #3}
\newcommand{\function}[5]{#1:#2\to #3,\quad #4\mapsto #5}
\newcommand{\pmqty}[1]{\begin{pmatrix}#1\end{pmatrix}}
\newcommand{\NN}{\mathbb{N}}
\newcommand{\ZZ}{\mathbb{Z}}
\newcommand{\QQ}{\mathbb{Q}}
\newcommand{\RR}{\mathbb{R}}
\newcommand{\degree}{^{\circ}}
\newcommand{\inverse}{^{-1}}
\newcommand{\compl}{^{\mathrm{c}}}
\newcommand{\qand}{\quad\text{and}\quad}
\newcommand{\qor}{\quad\text{or}\quad}
\newcommand{\qif}{\quad\text{if}\quad}
\newcommand{\qfor}{\quad\text{for}\quad}
\newcommand{\qqtext}[1]{\quad\text{#1}\quad}
\newcommand{\dd}[1]{\,\mathrm{d}#1}
\NewDocumentCommand{\dv}{s o m g}{\IfNoValueTF{#4}{\frac{\mathrm{d}}{\mathrm{d}#3}}{\IfNoValueTF{#2}{\frac{\mathrm{d}#3}{\mathrm{d}#4}}{\frac{\mathrm{d}^{#2}#3}{\mathrm{d}#4^{#2}}}}}
\newcommand{\SI}[2]{#1\,\text{#2}}
\newcommand{\centi}{c}
\newcommand{\metre}{m}
\newcommand{\percent}{\%}

\title{TMUA Lecture 1 -- Basic Algebra}
\author{Eason Shao; revised edition}
\date{2026-09-28}
\begin{document}
\maketitle
\paragraph{Editorial scope and provenance.}
Revised from Eason Shao's TMUA lecture notes; language and exercise attributions are retained.
The supplied licence assigns CC BY-NC-SA 4.0 to the text and GNU GPL 3.0 to the source files
(\url{https://creativecommons.org/licenses/by-nc-sa/4.0/};
\url{https://www.gnu.org/licenses/gpl-3.0.html}).
Revision: 2026-09-28. Missing personal-package macros have been defined locally;
selected mathematical errors and missing hypotheses have been corrected.
This is not an exhaustive certification of every exercise or its historical attribution.
Exercise alternatives and attempted proofs are questions to assess, not asserted results.
Keep each exercise with its statement, alternatives and any editorial explanation when indexing.
Figures are replaced by textual mathematical descriptions. Syllabus remarks describe the
source notes, not a verified current examination specification.
Unless stated otherwise, variables are real and \(\NN=\{1,2,3,\ldots\}\).
Explicit indices starting at zero override this convention.

\tableofcontents


\section{Factorisation Formulae}\label{tmua-1-section-001}
\subsection{General}\label{tmua-1-subsection-001}
\[
    \pqty{x \pm y}^2 = x^2 \pm 2xy + y^2, \quad \pqty{x \pm y}^3 = x^3 \pm 3 x^2 y + 3 x y^2 \pm y^3.
\]

\[
    x^2 - y^2 = \pqty{x + y} \pqty{x - y}.
\]

\[
    x^3 \pm y^3 = \pqty{x \pm y} \pqty{x^2 \mp xy + y^2}.
\]

\begin{example}\label{tmua-1-unit-001}
    Factorise \(x^4 - y^4\).
\end{example}

\begin{example}\label{tmua-1-unit-002}
    Factorise \(x^4 + 4 y^4\).
\end{example}

\begin{problem}\label{tmua-1-unit-003}
    When can \(x^n + y^n\) be factorised into a product of a linear factor and some other polynomial? What about \(x^n - y^n\)?
\end{problem}

\subsection{Binomial Theorem(s)}\label{tmua-1-subsection-002}
The \textbf{binomial theorem} is a systematic way to expand expressions of the form \(\pqty{a + b}^n\) where \(n \in \NN\). With the notation
\[
    \binom{n}{k} = \frac{n!}{k! \pqty{n - k}!},
\]
we have
\[
    \pqty{a + b}^n = \sum_{k = 0}^{n} \binom{n}{k} a^k b^{n - k}.
\]

Written out in full, it reads
\[
    \pqty{x + y}^n = \binom{n}{0} x^n y^0 + \binom{n}{1} x^{n - 1} y^1 + \cdots + \binom{n}{n - 1} x^1 y^{n - 1} + \binom{n}{n} x^0 y^n.
\]

\begin{example}[TMUA 2016 P1 Q1]\label{tmua-1-unit-004}
    It is given that the expansion of \(\pqty{ax + b}^3\) is \(8 x^3 - p x^2 + 18 x - 3 \sqrt{3}\), where \(a\), \(b\) and \(p\) are real constants. What is the value of \(p\)?
\end{example}

\begin{example}[TMUA 2017 P1 Q13]\label{tmua-1-unit-005}
    In the expansion of \(\pqty{a + bx}^5\) the coefficient of \(x^4\) is \(8\) times the coefficient of \(x^2\).

    Given that \(a\) and \(b\) are non-zero positive integers, what is the smallest possible value of \(a + b\)?
\end{example}

It is worth noting that the binomial coefficient has special cases of
\[
    \binom{n}{0} = 1, \quad \binom{n}{1} = n, \quad \binom{n}{2} = \frac{n \pqty{n - 1}}{2};
\]
and there is a symmetry of
\[
    \binom{n}{k} = \binom{n}{n - k}.
\]  

\textbf{Newton's Binomial Theorem} extends this to the form \(\pqty{1 + x}^{\alpha}\) where \(\alpha \in \RR\) and \(\abs{x} < 1\). It states that
\[
    \pqty{1 + x}^{\alpha} = 1 + \alpha x + \frac{\alpha \pqty{\alpha - 1}}{2} x^2 + \frac{\alpha \pqty{\alpha - 1} \pqty{\alpha - 2}}{3!} x^3 + \cdots + \frac{\alpha \pqty{\alpha - 1} \cdots \pqty{\alpha - n + 1}}{n!} x^n + \cdots
\]
For a nonnegative integer exponent, the expansion terminates; otherwise it is generally an infinite series, valid for the stated domain.

\begin{example}[TMUA 2016 P1 Q19]\label{tmua-1-unit-006}
    The coefficient of \(x^3\) in the expansion of \(\pqty{1 + 2x + 3 x^2}^6\) is equal to twice the coefficient of \(x^4\) in the expansion of \(\pqty{1 - a x^2}^5\). Find all possible values of the constant \(a\).
\end{example}

\begin{problem}\label{tmua-1-unit-007}
    Try proving the binomial theorem. Hint: prove and use the identity
    \[
        \binom{n + 1}{k} = \binom{n}{k} + \binom{n}{k - 1}.
    \]  
\end{problem}

\section{Exponents}\label{tmua-1-section-002}
The laws below hold for positive real bases and real exponents; in the root formula take a positive integer \(p\), and in reciprocal-exponent expressions take \(p\ne0\). Additional negative-base cases require separate domain checks.

\[
    a^b a^c = a^{b + c}, \quad \frac{a^b}{a^c} = a^{b - c};
\]
\[
    \pqty{a^b}^c = a^{bc}, \quad \pqty{ab}^p = a^p b^p;
\]
\[
    a^{-p} = \frac{1}{a^p}, \quad a^{\frac{1}{p}} = \sqrt[p]{a};
\]
\[
    a^0 = 1, \quad a^1 = a, \quad a^{-1} = \frac{1}{a}.
\]

\begin{example}[TMUA 2016 P1 Q11]\label{tmua-1-unit-008}
    The real roots of the equation \(4^{2x} + 12 = 2^{2x + 3}\) are \(p\) and \(q\), where \(p > q\). Find \(p - q\).
\end{example}

\begin{example}[TMUA 2017 P1 Q14]\label{tmua-1-unit-009}
    The solution of the simultaneous equations
    \begin{align*}
        2^x + 3 \times 2^y &= 3 \\
        2^{2x} - 9 \times 2^{2y} &= 6
    \end{align*}
    is \(\pqty{x, y} = \pqty{p, q}\). Find the value of \(p - q\).
\end{example}

\begin{problem}\label{tmua-1-unit-010}
    What are the restrictions to the \(a, b, c, p\) above?
\end{problem}

For \(a>0\), real powers are defined by \(a^t=\exp(t\log a)\). For negative bases, arbitrary real powers need not be real or defined; exponent laws cannot be extended blindly.

\section{Logarithms}\label{tmua-1-section-003}

For \(a > 0, a \neq 1\) and \(c > 0\), the \textbf{logarithm} base-\(a\) of \(c\) is defined as the unique \(b \in \RR\) such that \(a^b = c\).
\[
    a^b = c \iff b = \log_a c.
\]  

Laws of logarithms below assume \(a>0\), \(a\ne1\), \(x,y>0\), and \(k\ne0\) when dividing by \(k\):
\[
    \log_a x + \log_a y = \log_a \pqty{xy}, \quad \log_a x - \log_a y = \log_a \pqty{\frac{x}{y}};
\]
\[
    k \log_a x = \log_a \pqty{x^k}, \quad \frac{1}{k} \log_a x = \log_a \pqty{x^{\frac{1}{k}}};
\]
\[
    \log_a 1 = 0, \quad \log_a a = 1, \quad \log_a \frac{1}{a} = -1.
\]

\textbf{Change of base formula.} for \(a, b, c > 0, a \neq 1, c \neq 1\),
\[
    \log_a b = \frac{\log_c b}{\log_c a}.
\]

\begin{problem}\label{tmua-1-unit-011}
    Is there an equivalent version of change of base formula in exponentiation? How to prove the logarithm change of base formula?
\end{problem}

\begin{example}[TMUA 2016 P1 Q16]\label{tmua-1-unit-012}
    Given the simultaneous equations
    \begin{align*}
        \log_{10} 2 + \log_{10} \pqty{y - 1} &= 2 \log_{10} x \\
        \log_{10} \pqty{y + 3 - 3x} &= 0,
    \end{align*}
    find the values of \(y\).
\end{example}

\section{Surds}\label{tmua-1-section-004}

For \(a\in\RR\), \(b\ge0\), simplifying square roots: \(\sqrt{a^2 b} = \abs{a}\sqrt{b}\). \textbf{Do not forget to check whether the thing moved outside is positive or not.} Indeed, \((|a|\sqrt b)^2=a^2b\) and \(|a|\sqrt b\ge0\), so uniqueness of the nonnegative square root proves the identity.

Morale: \(\pqty{a + \sqrt{b}}\) multiplied with \(\pqty{a - \sqrt{b}}\) (its \textbf{conjugate surd}) gives \(a^2 - b\), which is rational when \(a,b\) are rational.

Application: rationalisation of denominators, assuming \(b\ge0\) and \(a^2\ne b\)
\[
    \frac{c}{a + \sqrt{b}} = \frac{c}{a + \sqrt{b}} \cdot \frac{a - \sqrt{b}}{a - \sqrt{b}} = \frac{c \pqty{a - \sqrt{b}}}{a^2 - b}.
\]

For \(a>0\):
\[
    \sqrt{\frac{1}{a}} = \frac{\sqrt{a}}{a}.
\]

The principal real square root of a nonnegative number is nonnegative.

\begin{example}\label{tmua-1-unit-013}
    Simplify the nested square roots \(\sqrt{49 \pm 12 \sqrt{5}}\) and \(\sqrt{7 \pm 2 \sqrt{10}}\).
\end{example}

\begin{example}[TMUA 2017 P2 Q8]\label{tmua-1-unit-014}
\textit{Exercise: the proposed statements or reasoning below must be assessed; they are not all asserted as true.}\par

    Which one of the following numbers is smallest in value?

    \begin{enumerate}
        \item \(\log_2 7\)
        \item \(\pqty{2^{-3} + 2^{-2}}^{-1}\)
        \item \(2^{\flatfrac{\pi}{3}}\)
        \item \(\dfrac{1}{4\pqty{\sqrt{2} - 1}^3}\)
        \item \(4 \sin^2 \pqty{\flatfrac{\pi}{4}}\)
    \end{enumerate}
\end{example}

\begin{problem}\label{tmua-1-unit-015}
    Find what is wrong with the following proof.

    \[
        1 = \sqrt{1} = \sqrt{\pqty{-1} \pqty{-1}} = \sqrt{-1} \sqrt{-1} = i^2 = -1.
    \]  
\end{problem}

\section{Quadratic Functions and Quadratic Equations}\label{tmua-1-section-005}
\subsection{Quadratic Formula}\label{tmua-1-subsection-003}

Consider the equation in \(x\): \(a x^2 + bx + c = 0\) (\(a \neq 0\)).

The \textbf{quadratic formula} reads:
\[
    x = \frac{-b \pm \sqrt{b^2 - 4ac}}{2a}
\]
where \(\Delta = b^2 - 4ac\) is the \textbf{discriminant}.

Proof is by completing the square.

\begin{align*}
    a x^2 + bx + c &= a \pqty{x^2 + \frac{b}{a} x + \frac{b^2}{4 a^2}} + c - \frac{b^2}{4a} \\
    &= a \pqty{x + \frac{b}{2a}}^2 + \frac{4ac - b^2}{4a}.
\end{align*}

This means the vertex of the graph (the optimum) has coordinates
\[
    \pqty{- \frac{b}{2a}, \frac{4ac - b^2}{4a}}.
\]

\subsection{Discriminant}\label{tmua-1-subsection-004}
The \textbf{discriminant} tells us about the roots of the equation.
\begin{itemize}
    \item If \(\Delta > 0\), then the equation has two distinct real roots.
    \item If \(\Delta = 0\), then the equation has two repeated (equal) real roots.
    \item If \(\Delta < 0\), then the equation has no real roots. (The roots are complex conjugates.)
\end{itemize}

The discriminant (together with the leading factor) also tells us whether the graph of the quadratic lies above/below or intersects the \(x\)-axis.

\begin{example}[TMUA 2017 P2 Q14]\label{tmua-1-unit-016}
    The diagram in \hyperref[fig:tmua-2017-p2-q14]{the textual figure description} shows the graph of \(y = x^2 - 2bx + c\). The vertex of this graph is at the point \(P\).

    Which one of the following could be the graph of \(y = x^2 - 2Bx + c\), where \(B > b\)?
\end{example}

\paragraph{Textual figure description.}\phantomsection\label{fig:tmua-2017-p2-q14}
Textual transcription of the graph choices. The original vertex is \(P=(b,c-b^2)\), with \(b>0\) and \(c-b^2>0\). Relative to \(P\), the candidate vertices are: A, right at the same height; B, left at the same height; C, directly above; D, directly below; E, above-right; F, below-right; G, above-left; H, below-left. All candidates are upward-opening parabolas.


\subsection{Vieta's Theorem}\label{tmua-1-subsection-005}
\textbf{Vieta's Theorem} states the two roots satisfy
\[
    x_1 + x_2 = - \frac{b}{a}, \quad x_1 x_2 = \frac{c}{a}.
\]

Proof is by quadratic formula (very disgusting), or by factorisation. We have
\[
    a x^2 + bx + c = a \pqty{x - x_1} \pqty{x - x_2}
\]
and expanding the brackets gives us the result.

\begin{example}\label{tmua-1-unit-017}
    Let \(x_1, x_2\) be the two real roots for the equation \(x^2 - 4x + 1 = 0\). Find:
    \begin{itemize}
        \item \(x_1^2 + x_2^2\);
        \item \(x_1^3 + x_2^3\);
        \item \(x_1^4 + x_2^4\).
    \end{itemize}
\end{example}

\begin{example}\label{tmua-1-unit-018}
    \textbf{Consistency check (corrected exercise).} Can the equation
\(x^2-(k-2)x+(k^2+3k+10)=0\) have real roots for any real \(k\)?
\textbf{Answer and justification.} No: its discriminant is
\[
(k-2)^2-4(k^2+3k+10)=-3(k+8/3)^2-44/3<0.
\]
Thus there is no maximum of the sum of squares of two real roots under these hypotheses.
\end{example}

\begin{example}[TMUA 2016 P2 Q15]\label{tmua-1-unit-019}
\textit{Exercise: the proposed statements or reasoning below must be assessed; they are not all asserted as true.}\par

    Consider the quadratic \(f\pqty{x} = x^2 - 2px + q\) and the statement \(f\pqty{x} = 0\) has two real roots whose difference is greater than \(2\) and less than \(4\).

    Which one of the following statements is true if and only if the above statement is true?

    \begin{enumerate}
        \item \(q < p^2 < q + 4\)
        \item \(\sqrt{q + 1} < p < \sqrt{q + 4}\)
        \item \(q - 3 \leq p^2 - 4 \leq q\)
        \item \(q < p^2 - 1 < q + 3\)
        \item \(q - 2 < p^2 - 3 < q + 2\)
    \end{enumerate}
\end{example}

\subsection{Alternative Forms of Quadratics}\label{tmua-1-subsection-006}
It is often more useful to look at a quadratic using its alternate forms:
\begin{itemize}
    \item \(a \pqty{x - p}^2 + q\) which tells us about its vertex \(\pqty{p, q}\);
    \item \(a \pqty{x - x_1} \pqty{x - x_2}\) which tells us about its roots \(x_{1, 2}\).
\end{itemize}

\section{Simultaneous Equations}\label{tmua-1-section-006}

Key strategies: change of variable, elimination, reduction of power.

If you have a equation in \(x\) and \(y\) and you want to find \(y\), express \(x\) in terms of \(y\), where possible, and substitute into the other equation (this is called \textbf{elimination}).

If the equation relates multiple powers of \(x\), then substituting other occurrences of e.g. \(x^2\) with lower powers is called \textbf{reduction of power}.

\section{Inequalities}\label{tmua-1-section-007}
\subsection{Simple Inequalities}\label{tmua-1-subsection-007}
\[
    a > b, b > c \implies a > c; \quad a > b \implies a + c > b + c;
\]
\[
    a > b, c > 0 \implies ac > bc; \quad a > b, c < 0 \implies ac < bc;
\]
\[
    a > b, c > d \implies a + c > b + d; \quad a > b > 0, c > d > 0 \implies ac > bd.
\]

\begin{example}[TMUA 2017 P2 Q13]\label{tmua-1-unit-020}
\textit{Exercise: the proposed statements or reasoning below must be assessed; they are not all asserted as true.}\par

    The positive real numbers \(a \times 10^{-3}\), \(b \times 10^{-2}\) and \(c \times 10^{-1}\) are each in standard form, and
    \[
        \pqty{a \times 10^{-3}} + \pqty{b \times 10^{-2}} = \pqty{c \times 10^{-1}}.
    \]

    Which of the following statements must be true?
    \begin{enumerate}
        \item \(a > 9\)
        \item \(b > 9\)
        \item \(a < c\)
        \item \(b < c\)
    \end{enumerate}
\end{example}

\subsection{Quadratic Inequalities}\label{tmua-1-subsection-008}
Quadratic inequalities can be solved by considering the leading coefficient and the discriminant. The most useful way to demonstrate is to sketch a graph.

\begin{example}[TMUA 2017 P1 Q19]\label{tmua-1-unit-021}
    The set of solutions to the inequality \(x^2 + bx + c < 0\) is the interval \(p < x < q\), where \(b, c, p\) and \(q\) are real constants with \(c < 0\). In terms of \(p, q\) and \(c\), find the set of solutions to the inequality \(x^2 + bcx + c^3 < 0\).
\end{example}

\subsection{Factorisable Inequalities}\label{tmua-1-subsection-009}
More complicated inequalities of the form
\[
    f_1\pqty{x} f_2\pqty{x} \cdots f_m\pqty{x} > 0
\]
can be solved by considering the critical values of \(f_i\) (i.e. solving \(f_i\pqty{x} = 0\)) and considering the signs in the respecting intervals split up by such critical values.

\begin{example}\label{tmua-1-unit-022}
    Solve the inequality \(\pqty{x - 2} \pqty{x + 1} \pqty{x - 1} > 0\).
\end{example}

\begin{example}\label{tmua-1-unit-023}
    Solve the inequality \(\pqty{x - 2}^2 \pqty{x + 1} < 0\).
\end{example}

\begin{example}[TMUA 2016 P2 Q8]\label{tmua-1-unit-024}
\textit{Exercise: the proposed statements or reasoning below must be assessed; they are not all asserted as true.}\par

    A region is defined by the inequalities \(x + y > 6\) and \(x - y > -4\). Consider the three statements:
    \begin{enumerate}
        \item \(x > 1\)
        \item \(y > 5\)
        \item \(\pqty{x + y} \pqty{x - y} > -24\)
    \end{enumerate}


    Which of the above statements is/are true for every point in the region?
\end{example}

\begin{example}[TMUA 2017 P1 Q5]\label{tmua-1-unit-025}
\textit{Exercise: the proposed statements or reasoning below must be assessed; they are not all asserted as true.}\par

    \(S\) is the complete set of values of \(x\) which satisfy \textbf{both} the inequalities \(x^2 - 8x + 12 < 0\) \textbf{and} \(2x + 1 > 9\). It can also be represented as a single inequality. Which one of the following single inequalities represents the set \(S\)?

    \begin{enumerate}
        \item \(\pqty{x^2 - 8x + 12} \pqty{2x + 1} < 0\)
        \item \(\pqty{x^2 - 8x + 12} \pqty{2x + 1} > 0\)
        \item \(x^2 - 10 x + 24 < 0\)
        \item \(x^2 - 10 x + 24 > 0\)
        \item \(x^2 - 6 x + 8 < 0\)
        \item \(x^2 - 6 x + 8 > 0\)
        \item \(x < 2\)
        \item \(x > 6\)
    \end{enumerate}
\end{example}


\end{document}
